We now present countermeasures against the attacks above, in three
categories: \emph{post-publication path checks},
which preserve both the signature size and the tree construction;
\emph{modifications to the path-construction algorithm}, which preserve the
signature size but change how the path is produced; and \emph{tree-less
constructions}, which remove the mechanism entirely at a large cost in
signature size.


\paragraph{Post-publication path check.}

These run checks on the output of $\SeedTreeToPath$. A size check alone already
defeats fault models~1 and~3, whose paths have length $1$ and \chg{$|\tpath|+k$}
respectively \chg{($|\tpath|$ the honest length, $k \geq 1$ the number of leaked
leaves)}, and the expected size is computable from $\calH$ in linear time.
A stronger version rebuilds the tree with $\PathToSeedTree$ and compares the
recovered leaves against $\calH$ and the reference tree.

\noindent\textit{Verify-before-release, and why it is weaker than it looks.}
The most natural version of this countermeasure is the classical one: run the
verification algorithm on the signature before releasing it, and suppress the
output on mismatch. Against the attacks of Section~\ref{sec:app} this does
work in principle, because none of the three faulted paths deserialises to the
honest leaf set and the recomputed digest therefore differs from $d$.

It is weaker than it looks, for a structural reason. In the MEDS reference
code $\SeedTreeToPath$ and $\PathToSeedTree$ are two modes of one function
sharing the traversal loop of Listing~\ref{lst:path}, so the instruction the
attacker glitches during signing is \emph{the same instruction}, at the same
code offset, that the verification pass executes. An attacker who can place
one glitch can place two, and a ChipWhisperer emits both within a single
trace: the redundancy is nominal, since the check re-executes the faulted code
rather than checking it independently. \chg{The replayed glitch does not merely
break the check, it makes it pass: in reconstruction mode the same skip makes
the traversal consume the path exactly as the faulted construction produced
it. Under Model~1, for instance, $\PathToSeedTree$ with the same skip reads
the single published seed as the root, regenerates the whole tree, and returns
every revealed leaf correctly, so the recomputed digest equals $d$. For
Models~2 and~3 the second glitch must also hit the same iteration (the same
value of \texttt{id}).} This is why we prefer the decoupled construction below, and the point generalises --- any
implementation folding $\Publish$ and $\Recon$ into shared code inherits it.
We have not measured the double-glitch success rate, and doing so is the most
useful single extension of this work. Finally, note that such a countermeasure
must suppress the output entirely on detection rather than merely flag it:
as Section~\ref{sec:app} observes, the adversary needs an emitted corrupted
path, not a valid signature.



\paragraph{Decoupling path construction from the flag tree.}

Coupling flag-tree construction to path extraction is what makes the replay
above possible. Decoupling them into independent stages with a check after
each means a fault during flag-tree generation shows up as a mismatch during
path construction:

\begin{enumerate}
\item \chg{Recompute the flag tree $\calF$ independently from $\mathbf{c}$
(equivalently $\calH$), in a pass separate from the one that serialises the
path.}
\item \chg{Verify \emph{leaf agreement in both directions}: $\calF[i]=1$ if and
only if $i\in\calH$, for \emph{every} leaf $i$ --- not only for the leaves
currently flagged $1$.}
\item \chg{Verify \emph{internal consistency}: every internal node satisfies
the OR-propagation relation $\calF(v)=\calF(\ell)\vee\calF(r)$, so no single
corrupted node can agree with its two children.}
\item \chg{Reconstruct the node path from the validated $\calF$ and the
reference seed tree, and assert that the published set is exactly the boundary
$\calF$ determines: no node with $\calF=1$ is derivable from it, and it contains
no node outside that boundary.}
\end{enumerate}

\chg{All four steps are $\mathcal{O}(t)$, against a signing cost dominated by
matrix arithmetic over $\Fq$. We implemented all three countermeasures and
timed whole signatures on a Cortex-M4 and an Intel Core i7
(Appendix~\ref{app:timings}). The differences from the baseline, between
$-8.0\%$ and $+14.5\%$, are of the same order as the run-to-run noise of such
measurements (several rows are \emph{faster} than the baseline), so they show
only that the overhead is small.}

\chg{\noindent\textit{What the check detects, and what it does not.} A weaker
form of this countermeasure tests only the ancestors of leaves \emph{currently}
flagged $1$, and so misses the very fault it should stop: a hidden leaf
$i\in\calH$ whose own flag is forced $\calF[i]\colon 1\to 0$ while its ancestors
stay $1$ is no longer examined, and if it then reaches the path the final
assertion---which also inspects only nodes flagged $1$---does not fire. The
if-and-only-if test of step~2 closes this, as it checks every leaf against
$\calH$ in both directions; step~3 catches any single internal-flag flip, since
such a flip breaks OR-propagation at its parent or one of its children. The
independent recomputation of step~1 is what gives the checks something to
disagree with: a fault during path construction cannot silently satisfy a
validation computed in a separate pass from a separate copy of $\mathbf{c}$.
Two faults remain outside the model. A double fault that corrupts the
construction pass and the independent validation identically---the shared-code
replay that defeats verify-before-release above---is not caught by redundancy
alone and needs the decoupling to be a genuinely distinct computation. A fault
that alters $\mathbf{c}$ itself before both passes read it is a
$\Chal$/$\Decide$ fault (Section~\ref{sec:taxonomy}), not one on $\Publish$, and
is outside this countermeasure's scope.}

\paragraph{Abandoning the tree.}

A third option is architectural: drop the tree and publish the ephemeral
matrices $(\tilde{A}_i,\tilde{B}_i)$ of the revealed rounds explicitly,
adapting to MEDS the LESS countermeasure of~\cite{ZKFault}. No seed is ever
released, so Invariant~\ref{inv:hiding} is enforced by construction rather
than by code. The cost is prohibitive where $t$ is large: MEDS-9923 grows from
$9{,}896$ to $677{,}440$ bytes, and even the best case, MEDS-167717, grows by
$68\%$. Appendix~\ref{app:treeless} gives the construction and the full
comparison.

